	\chapter{Rigidity and nuclear objects}
	\label[chapter]{chap:Rigidity_And_Nuclear_Objects}

	This chapter develops the symmetric monoidal counterpart of compact assembly. Trace-class morphisms take the role of compact morphisms, and nuclear objects are built from compact objects along trace-class maps. We first discuss rigidity and self-duality, then construct the rigidification, study nuclear objects and examine the analytic examples. The interpolation results of \Cref{chap:Compact_Morphisms_AB6} will be used in the rigidification argument.

	\section{Rigid symmetric monoidal categories}
	\label[section]{sec:Rigid_Categories}

	% Source: Talk 3 (Max Glaser, 27.10.2026); Clausen, Lecture 3 (19.07.2023); literature (Gaitsgory--Rozenblyum, Ramzi, Efimov), proofs partly completed with generative AI (Claude)
	A rich source of dualizable categories are \emph{rigid} symmetric monoidal categories, introduced by Gaitsgory--Rozenblyum \cite[Chapter~1, Section~9]{GaitsgoryRozenbluym2017StudyDAG}. We have already met them in \Cref{def:Rigid_Category}; we now give a more detailed treatment.

	\begin{definition}[{\cite[Chapter~1, Definition~9.1.2]{GaitsgoryRozenbluym2017StudyDAG}}]
		\label[definition]{def:Rigid_Category_Detailed}
		Let $C \in \CAlg(\PrL_{\st})$ be a presentably symmetric monoidal stable category, with unit $\unit_C$ and multiplication $\mu\colon C \otimes C \to C$. Then $C$ is called \emph{rigid} if
		\begin{enumerate}
			\item the unit $\unit_C$ is compact in $C$, or equivalently, the unit functor $u\colon \Sp \to C$, $\S \mapsto \unit_C$, is strongly continuous;
			\item the right adjoint $\mu^R\colon C \to C \otimes C$ preserves colimits, and its canonical lax $C \otimes C$-linear structure is strict; equivalently, $\mu$ admits a right adjoint in $\Mod_{C \otimes C}(\PrL_{\st})$.
		\end{enumerate}
	\end{definition}

	\begin{remark}
		\label[remark]{rem:Equivalent_Rigidity_Definitions}
		The equivalence of the two formulations in (2) is a general fact about adjunctions internal to module categories: a right adjoint of a $C \otimes C$-linear functor inherits a lax $C \otimes C$-linear structure, and it is a right adjoint in $\Mod_{C \otimes C}(\PrL_{\st})$ if and only if it preserves colimits and this lax structure is strict. \Cref{def:Rigid_Category_Detailed} is the definition of Gaitsgory--Rozenblyum, and it agrees with Ramzi's definition \cite[Definitions~4.5 and~4.36]{Ramzi_Locally_Rigid} (Ramzi additionally requires $C$ to be dualizable, which is automatic by \Cref{prop:Rigid_Self_Dual_Detailed}). It is also equivalent to \Cref{def:Rigid_Category}, which asks that for every finite set $I$ the multiplication $\mu_I\colon C^{\otimes I} \to C$ admits a right adjoint in $\Mod_{C^{\otimes I}}(\PrL_{\st})$. Indeed, for $I = \emptyset$ the map $\mu_{\emptyset}$ is the unit functor $u\colon \Sp \to C$, and $\Mod_{\Sp}(\PrL_{\st}) = \PrL_{\st}$, so this case is condition (1); the case $|I| = 1$ is vacuous, and the case $|I| = 2$ is condition (2). Conversely, if (2) holds, then by induction $\mu_{I \sqcup \{*\}} = \mu \circ (\mu_I \otimes \id_C)$ is a composite of two functors admitting right adjoints in $\Mod_{C^{\otimes I} \otimes C}(\PrL_{\st})$: for $\mu_I \otimes \id_C$ this follows by tensoring the adjunction for $\mu_I$ with $C$, and for $\mu$ by restricting the $C \otimes C$-linear adjunction along the algebra map $\mu_I \otimes \id_C\colon C^{\otimes I} \otimes C \to C \otimes C$.
	\end{remark}

	\begin{proposition}[{\cite[Chapter~1, Section~9.2.1]{GaitsgoryRozenbluym2017StudyDAG}}]
		\label[proposition]{prop:Rigid_Self_Dual_Detailed}
		Let $C$ be rigid. Then $C$ is self-dual in $\PrL_{\st}$, with evaluation and coevaluation given by
		\[
			\ev\colon C \otimes C \xrightarrow{\mu} C \xrightarrow{\Gamma := \hom_C(\unit_C,-)} \Sp, \qquad \coev\colon \Sp \xrightarrow{u} C \xrightarrow{\mu^R} C \otimes C.
		\]
		In particular, $C$ is dualizable.
	\end{proposition}
	\prooffor{prop:Rigid_Self_Dual_Detailed}
	\begin{proof}
		Note first that both functors preserve colimits: $\Gamma$ preserves colimits since $\unit_C$ is compact (and $C$ is stable), and $\mu^R$ preserves colimits by assumption. Moreover, $\Gamma$ is right adjoint to $u$. Write $\eta := \mu^R(\unit_C) \in C \otimes C$, so that $\coev$ is the colimit-preserving functor $\Sp \to C \otimes C$ sending $\S$ to $\eta$. We verify the first triangle identity, i.e.\ that the composite
		\[
			C \simeq \Sp \otimes C \xrightarrow{\coev \otimes C} C \otimes C \otimes C \xrightarrow{C \otimes \ev} C \otimes \Sp \simeq C
		\]
		is equivalent to the identity; the second one follows by symmetry. Since $C \otimes \ev = (C \otimes \Gamma) \circ (C \otimes \mu)$, the composite sends $c \in C$ to
		\[
			(C \otimes \Gamma)\big((C \otimes \mu)(\eta \boxtimes c)\big) \simeq (C \otimes \Gamma)\big(\eta \cdot (\unit_C \boxtimes c)\big),
		\]
		where $\cdot$ denotes the multiplication in $C \otimes C$. Now we use the linearity of $\mu^R$: for $a \in C \otimes C$ and $x \in C$ we have $\mu^R(\mu(a) \otimes x) \simeq a \cdot \mu^R(x)$. Taking $a = \unit_C \boxtimes c$ and $x = \unit_C$ gives
		\[
			\eta \cdot (\unit_C \boxtimes c) = (\unit_C \boxtimes c) \cdot \mu^R(\unit_C) \simeq \mu^R(c).
		\]
		Hence the composite is $(C \otimes \Gamma) \circ \mu^R\colon C \to C$, naturally in $c$ (as all functors involved preserve colimits, it suffices to identify the functors, which the above does naturally in $c$). Since $C \otimes - $ is a functor of $(\infty,2)$-categories, $C \otimes \Gamma$ is right adjoint to $C \otimes u\colon C \to C \otimes C$, $x \mapsto x \boxtimes \unit_C$. Hence $(C \otimes \Gamma) \circ \mu^R$ is right adjoint to $\mu \circ (C \otimes u) \simeq \id_C$, and is therefore equivalent to the identity.
	\end{proof}

	The linearity of the right adjoint to multiplication is what makes the triangle identity work. Trace-class morphisms express the corresponding duality data for individual morphisms.

	Next, we introduce trace-class morphisms, following Clausen--Scholze \cite{Scholze_Analytic_Geometry}.

	\begin{definition}
		\label[definition]{def:Trace_Class}
		Let $C$ be a symmetric monoidal category. A morphism $f\colon x \to y$ in $C$ is called \emph{trace-class} if there exist an object $z \in C$ and maps $\alpha\colon \unit_C \to y \otimes z$ and $\beta\colon z \otimes x \to \unit_C$ such that $f$ is homotopic to the composite
		\[
			x \simeq \unit_C \otimes x \xrightarrow{\alpha \otimes x} y \otimes z \otimes x \xrightarrow{y \otimes \beta} y \otimes \unit_C \simeq y.
		\]
	\end{definition}

	\begin{lemma}
		\label[lemma]{lem:Identity_Trace_Class}
		Let $C$ be a symmetric monoidal category whose homotopy category is idempotent complete; for instance, $C$ stable and idempotent complete \cite[Lemma~1.2.4.6]{lurie2016HA}. Then an object $x \in C$ is dualizable if and only if $\id_x$ is trace-class.
	\end{lemma}
	\prooffor{lem:Identity_Trace_Class}
	\begin{proof}
		If $x$ is dualizable with dual $x^{\vee}$, then the coevaluation and evaluation are a witness for $\id_x$. Conversely, let $(z,\alpha,\beta)$ be a witness for $\id_x$, so that $\alpha\colon \unit \to x \otimes z$ and $\beta\colon z \otimes x \to \unit$ satisfy the first triangle identity $(x \otimes \beta)(\alpha \otimes x) \simeq \id_x$, but not necessarily the second one. Consider
		\[
			e := (\beta \otimes z)(z \otimes \alpha)\colon z \to z \otimes x \otimes z \to z.
		\]
		Inserting the first triangle identity in the middle of $e \circ e$ shows that $e \circ e \simeq e$, and similarly $\beta(e \otimes x) \simeq \beta$ and $(x \otimes e)\alpha \simeq \alpha$. As the homotopy category is idempotent complete, $e$ splits: $e \simeq sr$ with $r\colon z \to z'$, $s\colon z' \to z$ and $rs \simeq \id_{z'}$. Put $\alpha' := (x \otimes r)\alpha$ and $\beta' := \beta(s \otimes x)$. Then
		\[
			(x \otimes \beta')(\alpha' \otimes x) \simeq (x \otimes \beta)(x \otimes e \otimes x)(\alpha \otimes x) \simeq (x \otimes \beta)(\alpha \otimes x) \simeq \id_x
		\]
		and $(\beta' \otimes z')(z' \otimes \alpha') \simeq r e s \simeq rsrs \simeq \id_{z'}$. So $(\alpha',\beta')$ exhibits $z'$ as a dual of $x$; since dualizability can be checked in the homotopy category, $x$ is dualizable.
	\end{proof}

	\begin{remark}
		\label[remark]{rem:Trace_Class_Basics}
		\leavevmode
		\begin{enumerate}
			\item If $C$ is closed, write $x^{\vee} := \iHom(x,\unit_C)$ for the \emph{predual} of $x$. Then $f$ is trace-class if and only if the map $\unit_C \to \iHom(x,y)$ classifying $f$ factors through the canonical map $y \otimes x^{\vee} \to \iHom(x,y)$: given $(z,\alpha,\beta)$ as above, $\beta$ is adjoint to a map $z \to x^{\vee}$, and composing $\alpha$ with it gives the desired factorization; conversely, given a factorization $\unit_C \to y \otimes x^{\vee}$, take $z = x^{\vee}$ and $\beta$ the evaluation map.
			\item Trace-class morphisms form a two-sided ideal: if $f$ is trace-class with witness $(z,\alpha,\beta)$, then $hfk$ is trace-class with witness $(z,(h \otimes z)\alpha,\beta(z \otimes k))$.
			\item If $C$ is stable and idempotent complete (for instance presentable and stable), then the identity of $x$ is trace-class if and only if $x$ is dualizable, see \Cref{lem:Identity_Trace_Class}.
			\item A symmetric monoidal functor preserves trace-class morphisms, since it preserves the data $(z,\alpha,\beta)$.
		\end{enumerate}
	\end{remark}

	\begin{lemma}
		\label[lemma]{lem:Trace_Class_Compact}
		Let $C \in \CAlg(\PrL_{\st})$ be such that $\unit_C$ is compact. Then every trace-class morphism in $C$ is compact in the sense of \Cref{def:Compact_Morphism}.
	\end{lemma}
	\prooffor{lem:Trace_Class_Compact}
	\begin{proof}
		Let $f\colon x \to y$ be trace-class with witness $(z,\alpha,\beta)$. We construct a factorization of $f^*\colon \Map_C(y,\colim(-)) \to \Map_C(x,\colim(-))$ through $\Map_{\Ind(C)}(y(x),-)$, naturally in $W = \quotcolim_j w_j \in \Ind(C)$. Given $g\colon y \to \colim_j w_j$, consider
		\[
			(g \otimes z) \circ \alpha\colon \unit_C \to (\colim_j w_j) \otimes z \simeq \colim_j (w_j \otimes z).
		\]
		Since $\unit_C$ is compact, this defines a point of $\colim_j \Map_C(\unit_C,w_j \otimes z)$, and postcomposing with the natural map
		\[
			\colim_j \Map_C(\unit_C,w_j \otimes z) \to \colim_j \Map_C(x,w_j), \qquad \gamma \mapsto (w_j \otimes \beta) \circ (\gamma \otimes x),
		\]
		we obtain a point of $\colim_j \Map_C(x,w_j) = \Map_{\Ind(C)}(y(x),W)$. All these maps are natural in $W$, and the image of this point in $\Map_C(x,\colim_j w_j)$ is $(\colim_j w_j \otimes \beta) \circ ((g \otimes z)\alpha \otimes x) \simeq g \circ (y \otimes \beta) \circ (\alpha \otimes x) \simeq gf$.
	\end{proof}

	The converse holds for rigid categories; this is the main structural result about rigid categories. It was formulated in Clausen's lectures as the \emph{definition} of rigidity for dualizable symmetric monoidal categories.

	\begin{theorem}[{Ramzi \cite[Corollaries~2.6, 4.55 and~4.60, Example~2.9]{Ramzi_Locally_Rigid}; see also \cite[Propositions~1.6 and~1.7]{Efimov_Inverse_Limits}}]
		\label[theorem]{thm:Rigid_Trace_Class}
		Let $C \in \CAlg(\PrL_{\st})$. Then the following are equivalent:
		\begin{enumerate}
			\item $C$ is rigid;
			\item $\unit_C$ is compact, and $C$ is generated under colimits by objects of the form $\colim(x_0 \to x_1 \to x_2 \to \cdots)$ in which each map $x_n \to x_{n+1}$ is trace-class;
			\item $C$ is dualizable, and a morphism in $C$ is compact if and only if it is trace-class.
		\end{enumerate}
	\end{theorem}

	We will not prove this theorem. Note however that some implications are formal given what we know: (1) $\implies$ the first half of (3) is \Cref{prop:Rigid_Self_Dual_Detailed}, and (3) $\implies$ (2) follows from \Cref{lem:Compact_Morphisms_Yhat}(2), since $\id_{\unit_C}$ is trace-class and hence compact, and \Cref{rem:Omega_1_Compact_Compact_Morphisms}, which shows that a dualizable category is generated by sequential colimits along compact morphisms. The hard part is that in a rigid category every compact morphism is trace-class \cite[Corollary~4.55, Example~2.9 and Corollary~2.6]{Ramzi_Locally_Rigid}, and that (2) implies (1) \cite[Corollary~4.60]{Ramzi_Locally_Rigid}. Note that the dualizability hypothesis in (3) was not stated in the talk.

	\begin{corollary}
		\label[corollary]{cor:Rigid_Functors_Strongly_Continuous}
		Let $F\colon C \to D$ be a colimit-preserving symmetric monoidal functor between rigid categories. Then $F$ is strongly continuous. Consequently, the category $\mathrm{Rig}$ of rigid symmetric monoidal categories and colimit-preserving symmetric monoidal functors is a full subcategory of $\CAlg(\Pr^{\dual})$.
	\end{corollary}
	\prooffor{cor:Rigid_Functors_Strongly_Continuous}
	\begin{proof}
		By \Cref{thm:Rigid_Trace_Class}, the compact morphisms in $C$ are trace-class. Since $F$ is symmetric monoidal, it preserves trace-class morphisms (\Cref{rem:Trace_Class_Basics}(4)), and trace-class morphisms in $D$ are compact by \Cref{lem:Trace_Class_Compact}. So $F$ preserves compact morphisms, and hence is strongly continuous by \Cref{prop:Strongly_Continuous_Compact_Morphisms}.
	\end{proof}

	\begin{remark}
		\label[remark]{rem:Strong_Continuity_From_Rigid_Source}
		The argument shows more generally that a colimit-preserving symmetric monoidal functor $F\colon C \to D$ with $C$ rigid and $D$ dualizable with compact unit is strongly continuous. In fact, the dualizability hypothesis on $D$ can be removed \cite[Proposition~1.9]{Efimov_Inverse_Limits}.
	\end{remark}

	\begin{example}
		\label[example]{ex:Rigid_Compactly_Generated}
		Let $C \in \CAlg(\PrL_{\st})$ be compactly generated. Then $C$ is rigid if and only if $\unit_C$ is compact and every compact object of $C$ is dualizable, if and only if the compact objects of $C$ are precisely the dualizable objects \cite[Chapter~1, Lemma~9.1.5 and Corollary~9.1.7]{GaitsgoryRozenbluym2017StudyDAG}. Indeed, if $\unit_C$ is compact, then dualizable objects are compact (as $\Map(x,-) \simeq \Map(\unit_C,x^{\vee} \otimes -)$), and conversely if all compact objects are dualizable, then $C = \Ind(C^{\omega})$ is generated by objects whose identity is trace-class, so (2) of \Cref{thm:Rigid_Trace_Class} holds. Examples include $\D(R)$ for a commutative ring $R$, more generally $\D_{\mathrm{qc}}(X)$ for a qcqs scheme $X$, and $\Sp$ itself.
	\end{example}

	\begin{example}
		\label[example]{ex:Shv_Compact_Hausdorff_Rigid}
		Let $X$ be a compact Hausdorff space. Then $\Shv(X;\Sp)$ is rigid; this can be seen as a manifestation of Verdier duality, see \cite[Example~1.4]{Efimov_Inverse_Limits}. For instance, the unit (the constant sheaf) is compact since global sections $\Gamma(X;-)$ commute with filtered colimits for compact $X$ \cite[Corollary~7.3.4.12]{lurie2009HTT}. In particular, $\Shv(X;\Sp)$ is dualizable and self-dual. For every presentable stable category $D$ we have $\Shv(X;D) \simeq \Shv(X;\Sp) \otimes D$ \cite[Proposition~2.8.7]{Krause_Nikolaus_Putzstuck_Sheaves}, so $\Shv(X;D)$ is dualizable whenever $D$ is. We will study these categories in detail in \Cref{chap:Sheaves}.
	\end{example}

	\begin{example}
		\label[example]{ex:Locally_Rigid_Sheaf_Categories}
		If $X$ is merely locally compact Hausdorff and not compact, then $\Shv(X;\Sp)$ is not rigid, since the unit is not compact. It is still dualizable, and even \emph{locally rigid} in the sense of Ramzi \cite[Definition~4.5 and Example~4.10]{Ramzi_Locally_Rigid} (see also \cite[Definition~1.36 and Example~1.38]{Efimov_Inverse_Limits} and \cite[Example~4.2.9, Corollary~4.4.13]{Krause_Nikolaus_Putzstuck_Sheaves}): it satisfies condition (2) of \Cref{def:Rigid_Category_Detailed}. Indeed, for $j\colon X \hookrightarrow X^+$ the one-point compactification, $j_!$ identifies $\Shv(X;\Sp)$ with a smashing ideal in the rigid category $\Shv(X^+;\Sp)$.
	\end{example}

	\section{Rigidification}
	\label[section]{sec:Rigidification}

	% Source: Clausen, Lecture 3 (19.07.2023); written from the literature (Ramzi, Construction 4.79 and Theorem 4.81; Efimov, Proposition 1.23 and Remark 1.24).
	A symmetric monoidal functor preserves trace-class maps, and in a rigid category the compact maps are precisely the trace-class maps. This suggests looking for a rigid category over $C$ through which every symmetric monoidal colimit-preserving functor from a rigid category factors.

	\begin{definition}[Rigidification]
		\label[definition]{def:Rigidification}
		Let $C$ be a presentably symmetric monoidal stable category. Its \emph{rigidification} is a rigid category $C^{\rig}$ with a symmetric monoidal colimit-preserving functor $C^{\rig}\to C$ such that composition induces an equivalence
		\[
			\Map_{\mathrm{Rig}}(R,C^{\rig})\simeq\Map_{\CAlg(\PrL_{\st})}(R,C)
		\]
		for every rigid category $R$.
	\end{definition}

	To construct it, choose an uncountable regular cardinal $\kappa$ such that $C\in\CAlg(\PrL_{\st,\kappa})$. Write $y\colon C^\kappa\to\Ind(C^\kappa)$ for the Yoneda embedding. Given a diagram $F\colon\Q_{\geq 0}\to C^\kappa$ whose transition maps $F(a)\to F(b)$ are trace-class in $C$ for all $a<b$, form its telescope
	\[
		\colim_{a\in\Q_{\geq 0}} y(F(a))\quad\in\Ind(C^\kappa).
	\]
	The full subcategory generated under colimits by these telescopes is $C^{\rig}$. The functor to $C$ takes the displayed formal colimit to the actual colimit of $F$ \cite[Construction~4.79 and Theorem~4.81]{Ramzi_Locally_Rigid}.

	Here are the main steps of Ramzi's proof. The tensor product of two telescopes is again a telescope, since the diagonal in $\Q_{\geq 0}\times\Q_{\geq 0}$ is cofinal and tensor products preserve trace-class maps. The constant telescope on the unit gives the compact object $y(\unit_C)$. To prove rigidity, Ramzi studies maps between telescopes that factor through a stage of the target. A cut at a later rational index gives an intermediate telescope through which such a map factors again. Ramzi shows that these maps detect nonzero objects and applies his rigidity criterion \cite[Corollary~4.59 and proof of Theorem~4.81]{Ramzi_Locally_Rigid}.

	For the universal property, let $R$ be rigid. Ramzi lifts a symmetric monoidal colimit-preserving functor $R\to C$ along $\Ind(C^\kappa)\to C$ using the lifting theorem for dualizable sources. The cardinal condition for this lift is automatic: compact maps in $R$ are trace-class, so their images in $C$ are trace-class. Since $\unit_C$ is $\kappa$-compact, these images are $\kappa$-compact maps, which implies that the functor preserves $\kappa$-compact objects. He then uses rational trace-class presentations in $R$ to show that the lift lands in the subcategory generated by the displayed telescopes. The equivalence of mapping animae in \Cref{def:Rigidification} follows from the lifting equivalence. In particular the construction is independent of the choice of $\kappa$ \cite[Theorem~4.81]{Ramzi_Locally_Rigid}.

	When $\unit_C$ is compact, the functor $C^{\rig}\to C$ is fully faithful. Indeed, trace-class maps in $C$ are then compact by \Cref{lem:Trace_Class_Compact}. Efimov observes that telescopes of such maps lie in the image of the partially defined left adjoint to realization $\Ind(C^\kappa)\to C$. On this image realization is inverse to the left adjoint. If moreover $C$ is compactly generated, each compact map factors through a compact object. In this case Efimov identifies $C^{\rig}$ inside $C$ as the subcategory generated under colimits by the actual colimits of rational trace-class diagrams in $C^\omega$ \cite[Proposition~1.23 and Remark~1.24(i)]{Efimov_Inverse_Limits}.

	\begin{remark}
		\label[remark]{rem:Rigidification_Via_Idempotent_Ideals}
		Clausen gives the same construction in his third lecture \citep{Clausen_Efimov_K_Theory_Lecture_3} as a special case of a construction for a presentable stable category $C$ and a two-sided ideal $S$ of maps. He assumes that $S$ contains the zero maps, is closed under finite sums and shifts in both directions, and satisfies the following cofiber condition. Given composable maps of arrows $u\to v\to w$ in $C$, if both maps on sources and both maps on targets lie in $S$, then $\cofib(u)\to\cofib(w)$ lies in $S$.

		Clausen forms $(C,S)^{\dual}$ by taking formal rational telescopes with transitions in $S$ and closing under colimits. For a dualizable category $D$, strongly continuous functors $D\to(C,S)^{\dual}$ correspond to colimit-preserving functors $D\to C$ that take compact maps into $S$. For symmetric monoidal $C$ he takes $S$ to be the trace-class maps to obtain $C^{\rig}$. Rational indexing matters here. If every map in $S$ factors as two maps in $S$, then \Cref{lem:Interpolation} lets us refine any sequence of maps in $S$ to a rational diagram with the same colimit. Only this last simplification uses the factorization hypothesis.
	\end{remark}

	\section{Nuclear objects}
	\label[section]{sec:Nuclear_Modules}

	% Source: Talk 3 (Max Glaser, 27.10.2026); literature (Clausen--Scholze, Efimov, Ramzi, Andreychev), drafted with generative AI (Claude)
	In condensed mathematics one constructs categories of nuclear modules inside larger categories of complete modules. We first study the abstract construction of Clausen and Scholze \cite{Scholze_Analytic_Geometry}. Nuclear objects can be built from compact objects along trace-class maps. We will use this description to determine when the nuclear objects form a rigid category.

	Throughout this section, let $C \in \CAlg(\PrL_{\st})$ be compactly generated with compact unit. We write $x^{\vee} = \iHom(x,\unit_C)$ for the predual of $x$.

	\begin{definition}[{\cite[Definition~1.25]{Efimov_Inverse_Limits}; cf.\ \cite[Definition~8.5]{Clausen_Scholze_Complex_Geometry} and \cite[Definitions~13.10 and~13.12]{Scholze_Analytic_Geometry}}]
		\label[definition]{def:Nuclear}
		An object $M \in C$ is called \emph{nuclear} if every morphism $P \to M$ from a compact object $P \in C^{\omega}$ is trace-class. We denote by $\Nuc(C) \subseteq C$ the full subcategory of nuclear objects. An object $M$ is called \emph{basic nuclear} if $M \simeq \colim(P_0 \to P_1 \to P_2 \to \cdots)$ with each $P_n$ compact and each map $P_n \to P_{n+1}$ trace-class.
	\end{definition}

	Note that a compact object $P$ is nuclear if and only if $\id_P$ is trace-class, i.e.\ if and only if $P$ is dualizable (\Cref{rem:Trace_Class_Basics}(3)). So the nuclear objects should be thought of as the objects which are ``locally dualizable''.

	For ordinary vector spaces, a tensor $v\otimes\ell\in x\otimes P^\vee$ gives the linear map $p\mapsto\ell(p)v$. The image of $x\otimes P^\vee\to\Hom(P,x)$ consists of the finite-rank maps. A finite-rank map $f\colon x\to y$ lets us go part of the way back: write $f(v)=\sum_i\ell_i(v)y_i$. Then for \emph{any} $b\colon P\to x$, the composite $fb$ is represented by the tensor $\sum_i y_i\otimes(\ell_i\circ b)$.

	The same construction works in $C$. The comparison $x\otimes P^\vee\to\iHom(P,x)$ is adjoint to evaluation on $P$. A trace-class presentation of $f$ plays the role of the expression $\sum_i y_i\otimes\ell_i$. Even when the comparison at $x$ has no inverse, a trace-class transition $x\to y$ lets us go diagonally from maps at $x$ to tensors at $y$.

	\begin{lemma}
		\label[lemma]{lem:Trace_Class_Interleaving}
		Let $f\colon x\to y$ be trace-class and let $P\in C$. The naturality square for the tensor-to-internal-Hom comparison admits a diagonal making both triangles commute up to homotopy:
		\[
			\begin{tikzcd}[column sep=large]
				x\otimes P^{\vee} \rar["f\otimes P^{\vee}"] \dar & y\otimes P^{\vee} \dar \\
				\iHom(P,x) \rar["{\iHom(P,f)}"'] \arrow[ur,dashed] & \iHom(P,y)
			\end{tikzcd}
		\]
	\end{lemma}
	\prooffor{lem:Trace_Class_Interleaving}
	\begin{proof}
		Choose a tensor $t\colon\unit_C\to y\otimes x^\vee$ representing $f$ as in \Cref{rem:Trace_Class_Basics}(1). For the diagonal, tensor with $t$ and then compose a functional on $x$ with a map from $P$ to $x$:
		\[
			\iHom(P,x)\xrightarrow{t\otimes\id}
			y\otimes x^\vee\otimes\iHom(P,x)\longrightarrow y\otimes P^\vee.
		\]
		The last arrow uses internal composition $x^\vee\otimes\iHom(P,x)\to P^\vee$, since $x^\vee=\iHom(x,\unit_C)$.

		For the lower triangle, evaluate on $P$. By associativity of composition, composing a functional with a map from $P$ and then evaluating on $P$ is the same as first evaluating the map into $x$ and then applying the functional. Contracting $t$ against this $x$ gives $f$, so the composite is postcomposition with $f$.

		For the upper triangle, start with $x\otimes P^\vee$. The composite $x^\vee\otimes x\otimes P^\vee\to x^\vee\otimes\iHom(P,x)\to P^\vee$ uses the evaluation $x^\vee\otimes x\to\unit_C$ and leaves the factor $P^\vee$ untouched. The composite around this triangle is therefore the contraction of $t$ against $x$, tensored with $\id_{P^\vee}$. This is $f\otimes P^\vee$.
	\end{proof}

	\begin{lemma}
		\label[lemma]{lem:Nuclear_Colimit_Interleaving}
		Let $M \simeq \colim_{i \in I} M_i$ for a directed poset $I$ such that for every $i$ there is $j \geq i$ with $M_i \to M_j$ trace-class. Let $P, Q \in C$ be such that $Q$ and $Q \otimes P$ are compact. Then the map $\hom_C(Q, M \otimes P^{\vee}) \to \hom_C(Q \otimes P, M)$ induced by $M \otimes P^{\vee} \to \iHom(P,M)$ is an equivalence of spectra.
	\end{lemma}
	\prooffor{lem:Nuclear_Colimit_Interleaving}
	\begin{proof}
		Apply $\pi_0\Map_C(Q,-)$ to the square of \Cref{lem:Trace_Class_Interleaving}. By adjunction the bottom row now consists of homotopy classes of maps out of $Q\otimes P$. Since $Q$ and $Q\otimes P$ are compact and tensoring with $P^\vee$ preserves colimits, the comparison at $M$ is, on $\pi_0$, the filtered colimit of the vertical comparisons at the stages $M_i$.

		For surjectivity, represent a class by a map $Q\otimes P\to M_i$. Choose a trace-class transition $M_i\to M_j$. The diagonal gives a class in $\pi_0\Map_C(Q,M_j\otimes P^\vee)$, and the lower triangle says it represents the required map at stage $j$.

		For injectivity, take two tensor classes with the same image in the colimit. Pass to a common stage where their images in the bottom row agree, and then choose a trace-class transition out of that stage. Applying the diagonal to the common image gives the same tensor class for both. By the upper triangle this is exactly the image of each original class at the later stage.

		This proves bijectivity on $\pi_0$. Replacing $Q$ by $\Sigma^n Q$ gives the same conclusion on $\pi_n$ for every $n\in\Z$, hence the equivalence of spectra.
	\end{proof}

	\begin{proposition}[Clausen--Scholze; see {\cite[Theorem~8.6]{Clausen_Scholze_Complex_Geometry} and \cite[Propositions~1.27 and~1.28]{Efimov_Inverse_Limits}}]
		\label[proposition]{prop:Nuclear_Basics}
		\leavevmode
		\begin{enumerate}
			\item For $M \in C$, the following are equivalent:
			\begin{enumerate}
				\item[(a)] $M$ is nuclear;
				\item[(b)] for every compact $P$, the map $\hom_C(\unit_C,P^{\vee} \otimes M) \to \hom_C(P,M)$ is an equivalence of spectra (equivalently, since we may replace $P$ by its shifts, the corresponding map of mapping spaces is an equivalence for every compact $P$);
				\item[(d)] $M \simeq \colim_{i \in I} M_i$ for a directed poset $I$ and compact objects $M_i$, such that for every $i$ there is $j \geq i$ with $M_i \to M_j$ trace-class.
			\end{enumerate}
			Moreover, these conditions are implied by
			\begin{enumerate}
				\item[(c)] for every compact $P$, the map $P^{\vee} \otimes M \to \iHom(P,M)$ is an equivalence,
			\end{enumerate}
			and if the tensor product of any two compact objects of $C$ is compact, then (c) is equivalent to (a), (b) and (d).
			\item The subcategory $\Nuc(C) \subseteq C$ is closed under colimits. It is $\omega_1$-presentable, and its $\omega_1$-compact objects are precisely the basic nuclear objects.
		\end{enumerate}
	\end{proposition}
	The main step is to turn the pointwise trace-class condition into a presentation whose transition maps are eventually trace-class. The interleaving lemma then upgrades surjectivity on $\pi_0$ to the isomorphism of mapping spectra in (b). For (2), we decompose these presentations into countable pieces, as in the proof of $\omega_1$-presentability for compactly assembled categories.
	\prooffor{prop:Nuclear_Basics}
	\begin{proof}
		We first prove (1). (b) $\Rightarrow$ (a): On $\pi_0$, surjectivity of the comparison in (b) says that every map $P\to M$ is represented by a tensor $\unit_C\to M\otimes P^\vee$. This is precisely the trace-class condition.

		(a) $\Rightarrow$ (d): Since $C$ is compactly generated, $M$ is a filtered colimit of compact objects, and by \Cref{thm:Filtered_Posets_Cofinal} we may write $M \simeq \colim_{i \in I} M_i$ with $I$ a directed poset and $M_i$ compact; denote by $t_{ij}$ the transition maps and by $\iota_i\colon M_i \to M$ the canonical maps. Fix $i$. Since $\iota_i$ is trace-class, it is represented by a tensor $w \in \pi_0\Map_C(\unit_C, M \otimes M_i^{\vee})$. As $\unit_C$ is compact and $M \otimes M_i^{\vee} \simeq \colim_j M_j \otimes M_i^{\vee}$, the class $w$ is the image of some $w' \in \pi_0\Map_C(\unit_C, M_j \otimes M_i^{\vee})$ with $j \geq i$. Let $g\colon M_i\to M_j$ be the map represented by $w'$. It is trace-class, and $\iota_jg\simeq\iota_i$ because the image of $w'$ at $M$ represents $\iota_i$. As $M_i$ is compact, $\pi_0\Map_C(M_i,M) \simeq \colim_k \pi_0\Map_C(M_i,M_k)$, so there is $k \geq j$ with $t_{jk}g \simeq t_{ik}$. Since trace-class morphisms form an ideal (\Cref{rem:Trace_Class_Basics}(2)), $t_{ik}$ is trace-class.

		(d) $\Rightarrow$ (b): This is \Cref{lem:Nuclear_Colimit_Interleaving} with $Q = \unit_C$.

		(c) $\Rightarrow$ (b): Apply $\hom_C(\unit_C,-)$ and use $\hom_C(\unit_C,\iHom(P,M)) \simeq \hom_C(P,M)$.

		(d) $\Rightarrow$ (c), if compact objects are closed under tensor products: For compact $Q$, the map $\hom_C(Q,P^{\vee} \otimes M) \to \hom_C(Q,\iHom(P,M)) \simeq \hom_C(Q \otimes P,M)$ is an equivalence by \Cref{lem:Nuclear_Colimit_Interleaving}. Since $C$ is compactly generated, the compact objects detect equivalences, so $P^{\vee} \otimes M \to \iHom(P,M)$ is an equivalence.

		We now prove (2). By (1), $M$ is nuclear if and only if the maps in (b) are equivalences. Since $\unit_C$ and $P$ are compact, the functors $\hom_C(\unit_C,P^{\vee} \otimes -)$ and $\hom_C(P,-)$ preserve filtered colimits, and being exact they preserve all small colimits. Hence $\Nuc(C)$ is closed under colimits in $C$; in particular, colimits in $\Nuc(C)$ are computed in $C$.

		Basic nuclear objects are nuclear by (d) (with $I = \N$). They are countable colimits of compact objects, hence $\omega_1$-compact in $C$ (\Cref{cor:Kappa_Small_Colimits_Of_Kappa_Compacts}), and hence also in $\Nuc(C)$.

		Next, every nuclear $M$ is an $\omega_1$-filtered colimit of basic nuclear objects. Write $M \simeq \colim_{i \in I} M_i$ as in (d), and choose for every $i$ an element $\sigma(i) \geq i$ such that $M_i \to M_{\sigma(i)}$ is trace-class. Let $\Sigma_{\sigma}$ be the set of countable directed subsets $J \subseteq I$ with $\sigma(J) \subseteq J$, ordered by inclusion. The proof of \Cref{prop:Colimits_Over_Countable_Subposets} applies verbatim to $\Sigma_{\sigma}$ instead of $\Sigma$ (adding in the inductive construction also $\sigma(u)$ for $u \in U_k$), and shows that $\Sigma_{\sigma}$ is $\omega_1$-filtered and that $M \simeq \colim_{J \in \Sigma_{\sigma}} M_J$ with $M_J = \colim_{i \in J} M_i$. Each $M_J$ is basic nuclear: enumerate $J = \{q_0,q_1,\dots\}$, and choose inductively $j_0 = q_0$ and $j_{n+1} \in J$ with $j_{n+1} \geq \sigma(j_n), q_{n+1}$. Then $n \mapsto j_n$ is cofinal (\Cref{lem:Cofinality_Posets}), so $M_J \simeq \colim_n M_{j_n}$, and each $M_{j_n} \to M_{j_{n+1}}$ factors through the trace-class map $M_{j_n} \to M_{\sigma(j_n)}$, hence is trace-class.

		The basic nuclear objects form an essentially small collection, since $C^{\omega}$ is essentially small. Thus $\Nuc(C)$ admits small colimits and is generated under $\omega_1$-filtered colimits by an essentially small collection of $\omega_1$-compact objects, so it is $\omega_1$-accessible by \cite[Proposition~5.4.2.2]{lurie2009HTT}, and hence $\omega_1$-presentable.

		Finally, let $X \in \Nuc(C)$ be $\omega_1$-compact. Writing $X \simeq \colim_{J \in \Sigma_{\sigma}} X_J$ as above, the identity of $X$ factors through some $X_J$, so $X$ is a retract of a basic nuclear object $Y = \colim_n a_n$, where the $a_n$ are compact and the transition maps $a_n \to a_{n+1}$ are trace-class. Now we argue as in the proof of \Cref{prop:Omega_1_Compact_In_Ind}, applied in $C \simeq \Ind(C^{\omega})$: by \Cref{lem:Retract_As_Telescope} and \Cref{lem:Strictification_Sequential_Maps}, we find a strictly increasing $f\colon \N \to \N$ with $f(n) \geq n$ and a natural transformation $\eta\colon a_{\bullet} \to a_{f(\bullet)}$ such that $X \simeq \colim_n \Phi(n,n)$, where $\Phi(k,n) = a_{f^k(n)}$. The transition map $\Phi(n,n) \to \Phi(n+1,n+1)$ is the composite of the transition map $a_{f^n(n)} \to a_{f^n(n+1)}$ of $a_{\bullet}$, which is trace-class as $f^n(n+1) > f^n(n)$, with $\eta_{f^n(n+1)}$. So it is trace-class, and $X$ is basic nuclear.
	\end{proof}

	\begin{remark}
		\label[remark]{rem:Nuclear_Definition_Pi0}
		\Cref{prop:Nuclear_Basics} shows that the condition in \Cref{def:Nuclear}, which only asks that the maps $\pi_0\Map_C(\unit_C,P^{\vee} \otimes M) \to \pi_0\Map_C(P,M)$ be surjective, already implies that $\hom_C(\unit_C,P^{\vee} \otimes M) \to \hom_C(P,M)$ is an equivalence. Clausen--Scholze \cite[Definition~8.5]{Clausen_Scholze_Complex_Geometry} instead define nuclearity by the latter condition (and \cite[Definition~13.10]{Scholze_Analytic_Geometry} by an analogous condition for analytic rings), whereas \Cref{def:Nuclear} follows \cite[Definition~1.25]{Efimov_Inverse_Limits}. The two definitions agree; the injectivity comes from the fact that, by (a) $\Rightarrow$ (d), $M$ is a filtered colimit of compact objects along trace-class maps, for which the two sides become interleaved (\Cref{lem:Nuclear_Colimit_Interleaving}).
	\end{remark}

	In the setting of this section, \Cref{sec:Rigidification} realizes $C^{\rig}$ inside $C$ by rationally indexed trace-class diagrams of compact objects. The following criterion lets us recognize the same objects using sequences.

	\begin{proposition}[{\cite[Proposition~1.29]{Efimov_Inverse_Limits}}]
		\label[proposition]{prop:Nuclear_Rigid}
		If every trace-class map between compact objects of $C$ factors through a compact object by two trace-class maps, then $C^{\rig}=\Nuc(C)$ as full subcategories of $C$.
	\end{proposition}
	\prooffor{prop:Nuclear_Rigid}
	\begin{proof}[Sketch]
		Each of the rationally indexed generators of $C^{\rig}$ from \Cref{sec:Rigidification} is nuclear by \Cref{prop:Nuclear_Basics}(1)(d). Since nuclear objects are closed under colimits, $C^{\rig}\subseteq\Nuc(C)$. Conversely, let $P_0 \to P_1 \to \cdots$ be a sequence of trace-class maps between compact objects. By the assumption and \Cref{lem:Interpolation}(2), applied to $D = C^{\omega}$ and the class of trace-class morphisms (which is closed under composition by \Cref{rem:Trace_Class_Basics}(2)), it extends to a functor $F\colon \Q_{\geq 0} \to C^{\omega}$ with trace-class transition maps and the same colimit. Restricting $F$ to the cofinal subset $\Q_{>0} \cong \Q$, we see that the basic nuclear objects lie in $C^{\rig}$. By \Cref{prop:Nuclear_Basics}(2), these generate $\Nuc(C)$.
	\end{proof}

	\section{Examples from condensed mathematics}
	\label[section]{sec:Condensed_Nuclear_Examples}

	The solid integers give a first example in which we can identify the nuclear objects explicitly. Passing to an adically complete ring leads to two categories of nuclear modules, obtained by rigidifying different ambient categories. We then look at individual nuclear modules over $\Z_p$ and at the corresponding questions over $\R$ and $\C$. In the solid and liquid examples, the compact objects are closed under tensor products, so the internal-Hom criterion of \Cref{prop:Nuclear_Basics}(1)(c) applies. We refer to \cite{Scholze_Analytic_Geometry} for the solid and liquid theory, to \cite{Clausen_Scholze_Complex_Geometry} for the gaseous theory, and to \cite{Andreychev_Thesis} and \cite[Section~7.2]{Efimov_Inverse_Limits} for the comparison with Efimov's theory.

	\begin{example}[Solid modules and nuclear modules over formal schemes]
		\label[example]{ex:Nuclear_Solid}
		Let $\D(\mathrm{Solid}_{\Z})$ be the derived category of light solid abelian groups, as in Efimov's convention in \cite[Section~7.2]{Efimov_Inverse_Limits}. This is a compactly generated symmetric monoidal category, with compact projective generator $\prod_{\N}\Z$ and compact unit $\Z$. It is far from rigid: the compact object $\prod_{\N}\Z$ is not dualizable (its predual is $\bigoplus_{\N}\Z$). In fact, there is a symmetric monoidal equivalence
		\[
			\D(\mathrm{Solid}_{\Z}) \simeq \Ind\big(\D^b(\mathrm{Proj}^{\omega_1}\text{-}\Z)\catop\big),
		\]
		where $\D^b(\mathrm{Proj}^{\omega_1}\text{-}\Z) \subseteq \D(\Z)$ is the stable subcategory generated by $\bigoplus_{\N}\Z$, under which $\prod_{\N}\Z$ corresponds to $(\bigoplus_{\N}\Z)\catop$ \cite[Section~7.2]{Efimov_Inverse_Limits}.

		In this example the nuclear objects admit a familiar description. A complex of light solid abelian groups is nuclear precisely when it comes from an ordinary complex of abelian groups via the discrete embedding \cite[Lecture~VIII, Exercise~1]{Clausen_Scholze_Complex_Geometry}. Thus $\Nuc(\D(\mathrm{Solid}_{\Z}))\simeq\D(\Z)$.

		More generally, let $R$ be a commutative ring and $I = (a_1,\dots,a_m) \subseteq R$ a finitely generated ideal, and let $R^{\wedge}_I$ denote its derived $I$-adic completion. We use the induced solid analytic structure, with the notation of \cite[Section~7.2]{Efimov_Inverse_Limits}. All completions and tensor products in this example are derived; for Noetherian $R$, the derived completion of $R$ agrees with its ordinary completion. Its category of modules $\D(\mathrm{Solid}_{(R^{\wedge}_I,\Z)}) = \Mod_{R^{\wedge}_I}(\D(\mathrm{Solid}_{\Z}))$ is compactly generated by $R^{\wedge}_I \otimes^{\square} \prod_{\N}\Z$, with compact unit. The category of \emph{nuclear solid modules} is
		\[
			\Nuc^{CS}(R^{\wedge}_I) := \Nuc\big(\D(\mathrm{Solid}_{(R^{\wedge}_I,\Z)})\big).
		\]
		Efimov identifies this category with the rigidification of the ambient solid module category:
		\[
			\Nuc^{CS}(R^{\wedge}_I) \simeq \D(\mathrm{Solid}_{(R^{\wedge}_I,\Z)})^{\rig}.
		\]

		The key technical input is that the preduals of compact objects are nuclear, which goes back to Andreychev \cite[Lemma~3.8 and Korollar~3.14]{Andreychev_Thesis}; see \cite[Propositions~7.4 and~7.5]{Efimov_Inverse_Limits} (Efimov deduces the rigidity from this via \cite[Propositions~1.33 and~1.34]{Efimov_Inverse_Limits} rather than via \Cref{prop:Nuclear_Rigid}). In particular:
		\begin{itemize}
			\item $\Nuc^{CS}(R^{\wedge}_I)$ is rigid, hence dualizable, but in general not compactly generated.
			\item Its compact objects are the perfect complexes over $R^{\wedge}_I$.
			\item Its continuous K-theory is $K^{\cont}(\Nuc^{CS}(R^{\wedge}_I)) \simeq \lim_n K(R \quot (a_1^n,\dots,a_m^n))$, which is $\lim_n K(R/I^n)$ if $R$ is Noetherian (Efimov \cite[Corollary~7.10]{Efimov_Inverse_Limits}, as conjectured by Clausen--Scholze), cf.\ \Cref{thm:Kcont_Nuclear}.
		\end{itemize}

		To compare the two constructions, keep the ring fixed and change the ambient category. Rigidifying the solid module category gives $\Nuc^{CS}(R^{\wedge}_I)$ as above. Rigidifying the category $\D_{I\text{-}\mathrm{compl}}(R)$ of derived $I$-complete $R$-modules gives Efimov’s variant $\Nuc(R^{\wedge}_I)$ \cite[Definition~3.28 and Proposition~4.3]{Efimov_Inverse_Limits}, which can be described as $\lim^{\dual}_n \D(R \quot (a_1^n,\dots,a_m^n))$, the limit computed in $\Pr^{\dual}$ \cite[Proposition~5.23]{Efimov_Inverse_Limits}. It receives a fully faithful strongly continuous functor from $\Nuc^{CS}(R^{\wedge}_I)$ \cite[Corollary~7.6]{Efimov_Inverse_Limits}, which induces an equivalence on compact objects, both sides having compact objects $\Perf(R^{\wedge}_I)$ \cite[Proposition~3.29]{Efimov_Inverse_Limits}. The ambient category of derived complete modules can itself fail to be rigid. For $R=\Z$ and $I=(p)$ its unit $\Z_p$ is not compact.

		% Source: written from the literature (Efimov, Localizing invariants of inverse limits, Corollary 7.20(i) and the discussion after Remark 7.21).
		Compact generation fails already for $\Nuc^{CS}(\Z_p)$. Indeed, suppose otherwise. Its compact objects are $\Perf(\Z_p)$, so the canonical $\Z$-linear functor $\D(\Z_p) \to \Nuc^{CS}(\Z_p)$ would be an equivalence. But Efimov computes
		\[
			\mathrm{HH}^{\cont}(\Nuc^{CS}(\Z_p)/\Z) \simeq \Z_p
		\]
		\cite[Corollary~7.20(i)]{Efimov_Inverse_Limits}, whereas
		\[
			\mathrm{HH}_1(\Z_p/\Z) \cong \Omega^1_{\Z_p/\Z} \cong \Omega^1_{\Q_p/\Q} \neq 0.
		\]
		Here $\mathrm{HH}^{\cont}$ denotes the continuous extension of relative Hochschild homology, which agrees with ordinary Hochschild homology on compactly generated categories. This gives the required contradiction; see also the discussion after \cite[Remark~7.21]{Efimov_Inverse_Limits}.

		% Source: written from the literature (Efimov, Corollary 7.7); moved from the introductory overview.
		There is also an operator-algebra presentation, analogous to the compact-operator example of \Cref{ex:Compact_Operators_Intro}. For the same pair $(R,I)$, take $P := (\bigoplus_{\N} R)^{\wedge}_I$ be the $I$-completed countable direct sum, which plays the role of a separable Hilbert space. In place of $B(\Hh)$ we take the endomorphism ring $\End_R(P)$, and in place of the Calkin algebra $B(\Hh)/\Kk$ the $I$-adically completed endomorphism ring
		\[
			\Calk^{\mathrm{top}}_R(P) := \End_{\Calk(R)}(P)^{\wedge}_I
		\]
		of $P$ in the Calkin category of $\Perf(R)$. The fiber of $\End_R(P) \to \Calk^{\mathrm{top}}_R(P)$ is $\big(\prod_{\N}R \otimes_R \bigoplus_{\N}R\big)^{\wedge}_I$, which one may think of as the $I$-adic completion of the endomorphisms of finite rank, just as $\Kk$ is the norm closure of the operators of finite rank. Efimov shows that $\End_R(P) \to \Calk^{\mathrm{top}}_R(P)$ is a homological epimorphism, and identifies the kernel of extension of scalars along this map with the category of \emph{nuclear} modules over $R^{\wedge}_I$ in the sense of Clausen--Scholze \cite[Corollary~7.7]{Efimov_Inverse_Limits}:
		\[
			\Nuc^{CS}(R^{\wedge}_I) \simeq \D\big(\End_R(P) \to \Calk^{\mathrm{top}}_R(P)\big).
		\]
	\end{example}

	\begin{example}[Banach spaces as nuclear modules]
		\label[example]{ex:Nuclear_Banach}
		Let $S=\N\cup\{\infty\}$ be the one-point compactification of the discrete set $\N$. A continuous function $S\to\Z_p$ is a convergent sequence together with its limit. Evaluation at $\infty$ splits by the constant functions, giving
		\[
			C(S,\Z_p)\simeq \Z_p\oplus c_0(\N,\Z_p),
		\]
		where $c_0(\N,\Z_p)$ consists of sequences tending to zero. By \cite[Lemma~3.8]{Andreychev_Thesis}, $C(S,\Z_p)$ is nuclear over $\Z_p^{\square}$ for every profinite $S$. Its summand $c_0(\N,\Z_p)$ is therefore nuclear, and inverting $p$ gives the nuclear solid $\Q_p$-vector space $c_0(\N,\Q_p)$.

		The convergence condition makes a difference. The module of all sequences, the compact objects such as $\prod_{\N}\Z_p$, which are of ``profinite'' or ``compact'' type. For instance, for every profinite set $S$ the module $C(S,\Z_p)$ of continuous functions is nuclear over $\Z_p^{\square}$ \cite[Lemma~3.8]{Andreychev_Thesis}. Taking $S = \N \cup \{\infty\}$, it follows that the $p$-adic Banach space $c_0(\N,\Z_p)$ of null sequences, a direct summand of $C(S,\Z_p)$, is nuclear, and inverting $p$ gives nuclear solid $\Q_p$-vector spaces such as $c_0(\N,\Q_p)$. In contrast, the compact object $\prod_{\N}\Z_p$ of $\D(\mathrm{Solid}_{\Z_p})$, is not nuclear: a compact nuclear object is dualizable, and the dualizable objects of $\D(\mathrm{Solid}_{\Z_p})$ are just the perfect complexes of $\Z_p$-modules \cite[Theorem~14.10]{Scholze_Analytic_Geometry}.
	\end{example}

	\begin{example}[Liquid and gaseous modules]
		\label[example]{ex:Nuclear_Liquid_Gaseous}
		Over the real and complex numbers, the solid theory is replaced by the \emph{liquid} theory: for $0 < p \leq 1$ there is a compactly generated symmetric monoidal category $\D(\mathrm{Liq}_p(\R))$ of $p$-liquid $\R$-vector spaces \cite{Scholze_Analytic_Geometry}, and one can form $\Nuc(\R^{p\text{-}\mathrm{liq}})$. The nuclear objects of $\D(\mathrm{Liq}_p)$ do not depend on $p$ \cite[Corollary~8.18]{Clausen_Scholze_Complex_Geometry}.

		According to Clausen \citep{Clausen_Efimov_K_Theory_Lecture_3}, there is an explicit model for the rigidification of the category of liquid $\C$-vector spaces in terms of operator algebras:
		\[
			\D^{\mathrm{liq}}(\C)^{\rig} \simeq \D\big(B(\Hh) \to B(\Hh)/S_{0+}\big).
		\]
		Here $\Hh$ is a separable infinite-dimensional complex Hilbert space and $B(\Hh)$ is its algebra of bounded operators. The ideal $S_{0+}=\bigcap_{q>0} S_q$ consists of compact operators whose singular values are $q$-summable for every $q>0$. As in the solid example, the category on the right is the kernel of extension of scalars along the displayed algebra map.

		Finally, the gaseous theory of Clausen--Scholze \cite{Clausen_Scholze_Complex_Geometry} provides yet another analytic ring structure on $\C$, with a category of nuclear modules $\Nuc(\C^{\mathrm{gas}})$. Rigidification produces rigid, hence dualizable, categories whose continuous K-theory is a candidate for analytic K-theory, cf.\ \Cref{sec:First_Computations}.
	\end{example}

	\begin{remark}
		\label[remark]{rem:Nuclear_Descent}
		On an affinoid adic space the categories of nuclear modules also glue. For a complete sheafy analytic Huber pair $(A,A^+)$, Andreychev proves that $U\mapsto\Nuc(\mathcal O(U))$ satisfies descent for covers by rational opens of $\operatorname{Spa}(A,A^+)$. The same holds after applying continuous K-theory, as recalled in \Cref{thm:Nuclear_Analytic_Descent}.
	\end{remark}
